\documentclass[10pt,a4paper]{article}
\usepackage{amsmath,amssymb,amsthm,graphicx}
\usepackage{hyperref}
\usepackage{cleveref}
\usepackage{amssymb,amsmath,amsthm}
\newcommand{\Mod}[1]{\ (\mathrm{mod}\ #1)}
\newcommand\blfootnote[1]{%
  \begingroup
  \renewcommand\thefootnote{}\footnote{#1}%
  \addtocounter{footnote}{-1}%
  \endgroup
}

\linespread{1.1}



\DeclareMathOperator{\Tr}{Tr}
\newtheorem{theorem}{Theorem}[section]
\newtheorem{lemma}{Lemma}[section]
\newtheorem{proposition}{Proposition}[section]
\newtheorem{corollary}{Corollary}[section]
\newtheorem{assumption}{Assumption}
\newtheorem{fact}{Fact}[section]
\newtheorem{claim}{Claim}
\newtheorem{question}{Question}
\newtheorem{remark}{Remark}
\newtheorem{definition}{Definition}[section]
\title{The spectral gap of a projection scan is at most its modified gap, which is at most the square root of n times the gap}
\author{Jason Gaitonde, Elchanan Mossel}
\date{arxiv:2410.11136v2; CC BY 4.0}
\begin{document}
\maketitle
\noindent\textit{Editorial note.} The counted statement is the author's labelled lemma environment at source0.tex lines 396-402; its proof environment opens at line 403, immediately after \textbackslash{}end\{lemma\}, so the pairing is by adjacency and the retained proof range 403-433 is the whole environment.\par\medskip
\noindent\textit{Editorial note.} This source is newly approved for this line and its preamble is built from the author's own package line, macro declarations and theorem-environment family (source lines 4-70); the heavy drawing, algorithm and font packages of the source are not retained because the retained material does not use them.\par\medskip
\noindent\textit{Editorial note.} Metadata provenance: the pinned source carries a licence observation file (CC BY 4.0, status supported, authors recorded there) but no separate source record file; per the project rule the missing metadata falls back to the licence observation, and the card records that here.\par\medskip
\noindent\textit{Editorial note.} The source carries no label-number dump and no counter values are imposed, so the printed numbering is the wrapper's default; the identity of the counted item is its author label lem:gibbs\_comp.\par\medskip
\noindent\textit{Editorial scope.} One original result of the pinned author manuscript, counted once: the two-sided comparison between the spectral gap, the modified gap and the coordinate gaps of a projection scan (author label lem:gibbs\_comp).\par\medskip
\noindent\textit{Allowed standard prerequisites (editorial).}\begin{itemize}
\item Projection scan. A projection scan is a product of transition matrices, one for each coordinate, applied in a fixed order; P = P\_n ... P\_1 is the transition matrix of the systematic scan.. It is the transition matrix whose gaps are compared.
\item Spectral gap and modified gap. The gap gamma of P is one minus the second largest eigenvalue in modulus, and the modified gap is the analogous quantity for the Dirichlet form of a single coordinate transition.. They are the two quantities appearing in the statement.
\item Dirichlet form. For a reversible chain with stationary measure pi, the Dirichlet form of a function is the average of the squared differences weighted by the edge conductances.. It is the quantity used to compare the two gaps.
\end{itemize}\medskip
\begin{lemma}
\label{lem:gibbs_comp}
    Let $\mathcal{X}=\mathcal{X}_1\times\ldots\times \mathcal{X}_n$ be a finite product space and let $\mathsf{P}=\mathsf{P}_n\ldots \mathsf{P}_1$ denote the transition matrix of the systematic scan with identity permutation.\footnote{Again, the actual Markov chain rerandomizes in the order $n,n-1,\ldots,1,n,\ldots$ by our convention.} Then
    \begin{equation*}
        \gamma_{\mathsf{P}}\leq \widetilde{\gamma_{\mathsf{P}}}\leq \sqrt{2n\gamma_{P}}.
    \end{equation*}
\end{lemma}

\begin{proof}
    The first inequality holds for any $\mathsf{P}$; indeed, for any $f$ such that $\mathbb{E}_{\pi}[f]=0$ and $\|f\|_{\pi}=1$, the reverse triangle inequality implies that
    \begin{equation*}
\|(I-\mathsf{P})f\|_{\pi}\geq \|f\|_{\pi}-\|Pf\|_{\pi}\geq 1-\|P\|_{\pi,\mathsf{op}}.
    \end{equation*}
    Since this holds for any such $f$, the first inequality follows.

    The second inequality is a special property of Gibbs samplers. Since each $\mathsf{P}_i$ is an orthogonal projection in $L^2(\mathcal{X},\pi)$, we will be able to precisely track the movement of $f$ under each projection. Let $f$ with $\mathbb{E}_{\pi}[f]=0$ and $\|f\|=1$ attain the maximum for $\|\mathsf{P}\|_{\pi,\mathsf{op}}$ so that
    \begin{equation*}
        (1-\gamma_{\mathsf{P}})^2=\|\mathsf{P}f\|_{\pi}^2
    \end{equation*} 
    Define $\gamma_0 = 0$ and for each $1\leq i\leq n$, set $\mathsf{Q}_i=\mathsf{P}_i\ldots \mathsf{P}_1$ and $\gamma_i$ via
    \begin{equation*}
        1-\gamma_i = \|\mathsf{Q}_if\|_{\pi}^2.
    \end{equation*}
    Note that this sequence is nondecreasing since norms are nonincreasing under projections and that $\gamma_n=2\gamma_{\mathsf{P}}-\gamma_{\mathsf{P}}^2\leq 2\gamma_{\mathsf{P}}$. Moreover, the Pythagorean theorem implies
    \begin{equation*}
        \|\mathsf{Q}_if-\mathsf{Q}_{i-1}f\|_{\pi}^2 = \|\mathsf{Q}_{i-1}f\|_{\pi}^2-\|\mathsf{Q}_if\|_{\pi}^2=\gamma_i-\gamma_{i-1}.
    \end{equation*}
    For this $f$, we thus have by the triangle inequality and Cauchy-Schwarz that
    \begin{align*}
        \|(I-\mathsf{P})f\|_{\pi}^2 &= \|f-\mathsf{Q}_nf\|_{\pi}^2\\
        &=\left\|\sum_{i=1}^n (\mathsf{Q}_if - \mathsf{Q}f)\right\|_{\pi}^2\\
        &\leq \left(\sum_{i=1}^n\left\|\mathsf{Q}_if - \mathsf{Q}_{i-1}f\right\|_{\pi}\right)^2\\
        &\leq n\cdot \sum_{i=1}^n\left\|\mathsf{Q}_if - \mathsf{Q}_{i-1}f\right\|_{\pi}^2\\
        &=n\cdot \sum_{i=1}^n(\gamma_i-\gamma_{i-1})\\
        &= n\cdot (\gamma_n-\gamma_0)\\
        &\leq 2n\gamma_{\mathsf{P}}.
    \end{align*}
Taking roots completes the proof by the variational formula of the second smallest singular value of the Laplacian.
\end{proof}

\end{document}
